\section{Space Complexity --- Exp BFS --- Main}
%%% generated by the blueprint pipeline from TCSlib.Complexity.SpaceComplexity.ExpBFS.Main
\subsection*{Space Complexity --- Exp BFS --- Main: Introduction}

Declarations of the Lean module \texttt{TCSlib.Complexity.SpaceComplexity.ExpBFS.Main}.

\subsection*{Space Complexity --- Exp BFS --- Main: Results}

\begin{definition}[Decider language family]
\lean{Complexity.LogProg.ExpBFS.decs}
\label{Complexity.LogProg.ExpBFS.decs}
\leanok
\uses{Complexity.Codec.acceptLang, Complexity.Codec.initLang, Complexity.Codec.stepLang, Complexity.inputLenLang, Turing.FinNDTM}
Given a finite nondeterministic Turing machine $N$ with binary alphabet, define a family
of four languages indexed by $j\in\{0,1,2,3\}$: in order, these are the input-length
language, the language of initial codes of $N$, the language of valid steps of $N$, and
the language of accepting codes of $N$.
\end{definition}

\begin{definition}[Language membership indicator]
\lean{Complexity.LogProg.ExpBFS.orcN}
\label{Complexity.LogProg.ExpBFS.orcN}
\leanok
\uses{Complexity.LogProg.ExpBFS.decs, Turing.FinNDTM, Turing.MultiTapeTM.indicator}
Given a finite nondeterministic Turing machine $N$, an index $j\in\{0,1,2,3\}$, and a
binary word $V$, define the Boolean value to be $1$ exactly when $V$ belongs to the
$j$th language in the four-language family associated with $N$, and $0$ otherwise.
\end{definition}

\begin{lemma}[Input-Length Oracle on Encoded Indices]
\lean{Complexity.LogProg.ExpBFS.orcN_len}
\label{Complexity.LogProg.ExpBFS.orcN_len}
\leanok
\uses{Complexity.LogProg.ExpBFS.decs, Complexity.LogProg.ExpBFS.orcN, Complexity.indicator_indexLang_pairEncode, Complexity.inputLenLang, Turing.FinNDTM, Turing.pairEncode}
\difficulty{3}
For every finite Boolean string $x$ and natural number $k\le |x|$, the first oracle
associated with the nondeterministic machine $N$, queried on the pair encoding of $x$
and the binary representation of $k$, returns true exactly when $k<|x|$.
\end{lemma}

\begin{lemma}[Initial Code Oracle Characterization]
\lean{Complexity.LogProg.ExpBFS.ini_orcN}
\label{Complexity.LogProg.ExpBFS.ini_orcN}
\leanok
\uses{Complexity.Codec.cfgCode, Complexity.Codec.mem_initLang, Complexity.LogProg.ExpBFS.decs, Complexity.LogProg.ExpBFS.ini, Complexity.LogProg.ExpBFS.orcN, Complexity.LogProg.indicator_eq_decide, Turing.NDTM.initCfg}
\difficulty{3}
Let $N$ be a finite nondeterministic Turing machine, let $x$ be a bit string, and let
$r,u\in\mathbb N$. The initial-code oracle on the paired input consisting of $x$ and the
nested binary encodings of $u$ and $r$ accepts exactly when $u$ is the configuration
code of $N$'s initial configuration on input $x$, using window parameter $r$.
\end{lemma}

\begin{lemma}[Step Oracle Characterization]
\lean{Complexity.LogProg.ExpBFS.stp_orcN}
\label{Complexity.LogProg.ExpBFS.stp_orcN}
\leanok
\uses{Complexity.Codec.StepRel, Complexity.LogProg.ExpBFS.decs, Complexity.LogProg.ExpBFS.orcN, Complexity.LogProg.ExpBFS.stp, Complexity.LogProg.indicator_eq_decide}
\difficulty{1}
Let $N$ be a finite nondeterministic Turing machine over the Boolean alphabet, let $x$
be a finite bit string, and let $r,u,v\in\mathbb N$. The step oracle associated with $N$
returns true on the nested encoding of $x$ and the binary representations of $r,u,v$ if
and only if the encoded pair of configurations with codes $u$ and $v$, at window size
$r$, belongs to the step language of $N$.
\end{lemma}

\begin{lemma}[Oracle acceptance characterization]
\lean{Complexity.LogProg.ExpBFS.acc_orcN}
\label{Complexity.LogProg.ExpBFS.acc_orcN}
\leanok
\uses{Complexity.Codec.acceptLang, Complexity.LogProg.ExpBFS.decs, Complexity.LogProg.ExpBFS.orcN, Complexity.LogProg.indicator_eq_decide, Turing.pairEncode}
\difficulty{1}
Let $N$ be a finite-state nondeterministic Turing machine over the Boolean alphabet, let
$x$ be a bit string, and let $r,u$ be natural numbers. The oracle for $N$, queried at
index $3$ on the encoding of $x$ paired with the binary encodings of $u$ and $r$,
returns true if and only if that encoded string belongs to the acceptance language of
$N$.
\end{lemma}

\begin{lemma}[Oracle reachability and machine configuration codes]
\lean{Complexity.LogProg.ExpBFS.cl_iff}
\label{Complexity.LogProg.ExpBFS.cl_iff}
\leanok
\uses{Complexity.Codec.cfgCode, Complexity.Codec.cfgCode_lt, Complexity.Codec.codeK, Complexity.ExpBFS.RCode, Complexity.ExpBFS.rcode_lt, Complexity.LogProg.ExpBFS.Cl, Complexity.LogProg.ExpBFS.ini_orcN, Complexity.LogProg.ExpBFS.orcN, Complexity.LogProg.ExpBFS.stp, Complexity.LogProg.ExpBFS.stp_orcN, Complexity.logSpace, Turing.NDTM.initCfg}
\difficulty{5}
Let $N$ be a nondeterministic Turing machine with a finite state set and Boolean
work-tape alphabet, let $x$ be a bit string, and let $r,K$ be natural numbers. If
$2^{(2W+2)(r+\mathrm{logSpace}(|x|)+1)}\le K$, where $W=3(|\mathrm{State}(N)|+1)+3k+1$
and $k$ is the number of work tapes of $N$, then a number $v$ is reachable from some
$u<K$ accepted by the initial-configuration oracle, through a finite chain of
step-oracle transitions whose targets are all less than $K$, if and only if $v$ is
reachable from the code of the initial configuration of $N$ on $x$ by a finite chain of
valid one-step configuration codes, with each configuration confined to work-tape
positions of absolute value at most $r$.
\end{lemma}

\begin{lemma}[Equality of bounded reachability counts]
\lean{Complexity.LogProg.ExpBFS.card_cl}
\label{Complexity.LogProg.ExpBFS.card_cl}
\leanok
\uses{Complexity.Codec.codeK, Complexity.ExpBFS.card_filter_rcode, Complexity.ExpBFS.rset, Complexity.LogProg.ExpBFS.Cl, Complexity.LogProg.ExpBFS.cl_iff, Complexity.LogProg.ExpBFS.orcN, Complexity.logSpace}
\difficulty{3}
Let $N$ be a nondeterministic Turing machine with a finite state set $S$ and $k$ work
tapes, using the Boolean alphabet, and let $x$ be a Boolean string. Put
$\ell=\lfloor\log_2 |x|\rfloor+1$, $W=3(|S|+1)+3k+1$, and $q=2W+2$. For natural numbers
$r$ and $K$ satisfying $2^{q(r+\ell+1)}\le K$, the number of integers $v<K$ reachable
from some $u<K$ accepted by the initial oracle, by a finite chain of steps accepted by
the step oracle whose successive targets are all less than $K$, equals the number of
distinct configuration codes reachable from the initial configuration code of $N$ on $x$
by machine steps, with codes restricted to be less than $2^{q(r+\ell+1)}$.
\end{lemma}

\begin{theorem}[Halting Characterisation of the Final Scan]
\lean{Complexity.LogProg.ExpBFS.fscan_run}
\label{Complexity.LogProg.ExpBFS.fscan_run}
\leanok
\uses{Complexity.LogProg.ExpBFS.A, Complexity.LogProg.ExpBFS.G, Complexity.LogProg.ExpBFS.SBd, Complexity.LogProg.ExpBFS.Site.idx, Complexity.LogProg.ExpBFS.Site.set, Complexity.LogProg.ExpBFS.St, Complexity.LogProg.ExpBFS.St.Bd, Complexity.LogProg.ExpBFS.St.toV_U, Complexity.LogProg.ExpBFS.acExit, Complexity.LogProg.ExpBFS.ac_run, Complexity.LogProg.ExpBFS.cell, Complexity.LogProg.ExpBFS.cf, Complexity.LogProg.ExpBFS.g_of, Complexity.LogProg.ExpBFS.stepH_of, Complexity.LogProg.MHalt, Complexity.LogProg.MHalt.ret, Complexity.LogProg.MReach.halt, Complexity.LogProg.astep, Complexity.LogProg.vword_whole₂, Turing.pairEncode}
\difficulty{6}
Let $K_0,B,P,j$ be natural numbers, let $x$ be a Boolean word, let $\mathcal O$ assign a
Boolean answer to each of four indices and each Boolean word, let
$S=(N,L_g,s,s_1,X,w,K,B_s,U,V,T,T_2,Z,C,ch,c,c_1)$ be a machine state, and let
$M:\mathbb Z\to\{\mathrm{false},\mathrm{true}\}$ be a memory marking. Suppose $K-U=j$
using truncated subtraction, $U\le K$, $T=T_2=Z=0$, every coordinate of $S$ is at most
$B$, $B_s+K\le B$, $B_s+K\le P$, and $K+1\le B$. Then the machine started in its
final-scan state with registers given by $S$, memory $M$, and memory head at $0$ halts
while its guard holds at every time strictly before halting. Its answer is true exactly
when there exists $v\in\mathbb N$ such that $U\le v<K$, $M(B_s+v)=\mathrm{true}$, and
$\mathcal O$ at index $3$ answers true on the encoding of $x$ paired with the encodings
of $v$ and $s$.
\end{theorem}

\begin{definition}[Exponential machine bound]
\lean{Complexity.LogProg.ExpBFS.kb}
\label{Complexity.LogProg.ExpBFS.kb}
\leanok
\uses{Complexity.Codec.codeK, Complexity.LogProg.ExpBFS.lgv, Turing.FinNDTM}
For a finite nondeterministic Turing machine $N$ and natural numbers $n$ and $s$, define
$K_N(n,s)=2^{c(N)(s+L(n)+1)}$, where $c(N)$ is the machine code parameter and $L(n)$ is
the logarithmic quantity used in the definition.
\end{definition}

\begin{lemma}[Lower bound for the code bound]
\lean{Complexity.LogProg.ExpBFS.two_le_codeK}
\label{Complexity.LogProg.ExpBFS.two_le_codeK}
\leanok
\uses{Complexity.Codec.codeK, Turing.FinNDTM}
\difficulty{2}
For every finite nondeterministic Turing machine $N$ over the Boolean alphabet, let $H$
be the number of elements of $(	ext{Option } 	ext{State}(N))\times	ext{Fin}(3)$ and set
$W=H+3k+1$, where $k$ is the number of work tapes. If its code bound is $2W+2$, then
this bound is at least $2$.
\end{lemma}

\begin{lemma}[Growth of the Exponential Bound]
\lean{Complexity.LogProg.ExpBFS.four_mul_kb_le}
\label{Complexity.LogProg.ExpBFS.four_mul_kb_le}
\leanok
\uses{Complexity.Codec.codeK, Complexity.LogProg.ExpBFS.kb, Complexity.LogProg.ExpBFS.lgv, Complexity.LogProg.ExpBFS.two_le_codeK, Turing.FinNDTM}
\difficulty{4}
For every finite nondeterministic Turing machine $N$ over the Boolean alphabet and all
natural numbers $n,s$, the bound $k_b(N,n,s)$ satisfies $4k_b(N,n,s)\le k_b(N,n,s+1)$.
\end{lemma}

\begin{lemma}[Monotonicity of the exponential code bound]
\lean{Complexity.LogProg.ExpBFS.kb_mono}
\label{Complexity.LogProg.ExpBFS.kb_mono}
\leanok
\uses{Complexity.LogProg.ExpBFS.kb, Turing.FinNDTM}
\difficulty{2}
For every finite-state nondeterministic Turing machine $N$, every $n\in\mathbb{N}$, and
natural numbers $s\le s'$, the exponential bound $kb(N,n,s)$ is at most $kb(N,n,s')$.
\end{lemma}

\begin{lemma}[Exponential code bound]
\lean{Complexity.LogProg.ExpBFS.code_bound_le}
\label{Complexity.LogProg.ExpBFS.code_bound_le}
\leanok
\uses{Complexity.Codec.codeK, Complexity.LogProg.ExpBFS.kb, Complexity.LogProg.ExpBFS.logSpace_le_lgv, Complexity.logSpace, Turing.FinNDTM}
\difficulty{3}
Let $N$ be a finite nondeterministic Turing machine over the Boolean alphabet, and let
$n,s\in\mathbb{N}$. Then $2^{c(N)(s+\log_2^+(n)+1)}\leq 2^{c(N)(s+g(n)+1)}$, where
$c(N)$ is the machine’s code-size parameter, $\log_2^+(n)=\lfloor\log_2 n\rfloor+1$
(with value $1$ at $n=0$), and $g(n)=1$ when $n=0$ and $g(n)=\lfloor\log_2 n\rfloor+2$
when $n>0$.
\end{lemma}

\begin{definition}[Buffer size bound]
\lean{Complexity.LogProg.ExpBFS.bigB}
\label{Complexity.LogProg.ExpBFS.bigB}
\leanok
\uses{Complexity.LogProg.ExpBFS.kb, Turing.FinNDTM}
For a finite nondeterministic Turing machine $N$ and natural numbers $n,t$, define the
buffer size bound as three times the search bound for parameters $n$ and $t+1$, plus
$n+3$.
\end{definition}

\begin{lemma}[Filtered initial interval bound]
\lean{Complexity.LogProg.ExpBFS.card_filter_range_le}
\label{Complexity.LogProg.ExpBFS.card_filter_range_le}
\leanok
\difficulty{1}
For every natural number $K$ and every decidable predicate $p$ on the natural numbers,
the number of elements of $[0,K)$ satisfying $p$ is at most $K$.
\end{lemma}

\begin{lemma}[Exponential bound for the encoded parameter]
\lean{Complexity.LogProg.ExpBFS.lt_kb}
\label{Complexity.LogProg.ExpBFS.lt_kb}
\leanok
\uses{Complexity.Codec.codeK, Complexity.LogProg.ExpBFS.kb, Complexity.LogProg.ExpBFS.lgv, Turing.FinNDTM}
\difficulty{1}
Let $N$ be a finite nondeterministic Turing machine over the Boolean alphabet, and let
$n,s$ be natural numbers. Set $c_N=2W_N+2$ and define $h(0)=0$, while for $n>0$ set
$h(n)=\lfloor\log_2 n\rfloor+1$. If $K=2^{c_N(s+h(n)+2)}$, then
$c_N(s+h(n)+2)<K$.
\end{lemma}

\begin{definition}[Initial search state]
\lean{Complexity.LogProg.ExpBFS.st1}
\label{Complexity.LogProg.ExpBFS.st1}
\leanok
\uses{Complexity.LogProg.ExpBFS.St, Complexity.LogProg.ExpBFS.St.zero, Complexity.LogProg.ExpBFS.lgv}
For natural numbers $n$, $s$, and $c_1$, let the initial search state be the state
record with all fields zero except that its logarithm-bound field is set to the
logarithm-bound value associated to $n$, its radius field to $s$, and its first-search
count field to $c_1$.
\end{definition}

\begin{definition}[Second search state]
\lean{Complexity.LogProg.ExpBFS.st2}
\label{Complexity.LogProg.ExpBFS.st2}
\leanok
\uses{Complexity.LogProg.ExpBFS.St, Complexity.LogProg.ExpBFS.St.zero, Complexity.LogProg.ExpBFS.lgv}
For natural numbers $n$, $s$, and $c_1$, let the second search state be the state record
with all fields zero except that its logarithm-bound field is set to the logarithm-bound
value associated to $n$, its radius and next-radius fields both to $s$, and its
first-search count field to $c_1$.
\end{definition}

\begin{definition}[Third search state]
\lean{Complexity.LogProg.ExpBFS.st3}
\label{Complexity.LogProg.ExpBFS.st3}
\leanok
\uses{Complexity.LogProg.ExpBFS.St, Complexity.LogProg.ExpBFS.St.zero, Complexity.LogProg.ExpBFS.lgv}
For natural numbers $n$, $s$, and $c_1$, let the third search state be the state record
with all fields zero except that its logarithm-bound field is set to the logarithm-bound
value associated to $n$, its radius field to $s$, its next-radius field to $s+1$, and
its first-search count field to $c_1$.
\end{definition}

\begin{definition}[Fifth search state]
\lean{Complexity.LogProg.ExpBFS.st5}
\label{Complexity.LogProg.ExpBFS.st5}
\leanok
\uses{Complexity.LogProg.ExpBFS.St, Complexity.LogProg.ExpBFS.St.zero, Complexity.LogProg.ExpBFS.lgv}
For natural numbers $n$, $s$, and $c_1$, let the fifth search state be the state record
with all fields zero except that its logarithm-bound field is set to the logarithm-bound
value associated to $n$, its radius field to $s$, its next-radius field to $s+1$, and
both its marked-code count and first-search count fields to $c_1$.
\end{definition}

\begin{theorem}[Bounded-space search returns the correct answer]
\lean{Complexity.LogProg.ExpBFS.iter_run}
\label{Complexity.LogProg.ExpBFS.iter_run}
\leanok
\uses{Complexity.Codec.codeK, Complexity.ExpBFS.card_rset_succ_of_le, Complexity.ExpBFS.mem_iff_of_card_eq, Complexity.ExpBFS.rcode_lt, Complexity.ExpBFS.rset, Complexity.LogProg.ExpBFS.A, Complexity.LogProg.ExpBFS.Cl, Complexity.LogProg.ExpBFS.G, Complexity.LogProg.ExpBFS.SBd, Complexity.LogProg.ExpBFS.St, Complexity.LogProg.ExpBFS.St.Bd, Complexity.LogProg.ExpBFS.St.zero, Complexity.LogProg.ExpBFS.acc_orcN, Complexity.LogProg.ExpBFS.bigB, Complexity.LogProg.ExpBFS.card_cl, Complexity.LogProg.ExpBFS.card_filter_range_le, Complexity.LogProg.ExpBFS.cell, Complexity.LogProg.ExpBFS.cf, Complexity.LogProg.ExpBFS.cl_iff, Complexity.LogProg.ExpBFS.code_bound_le, Complexity.LogProg.ExpBFS.cpDst, Complexity.LogProg.ExpBFS.cpExit, Complexity.LogProg.ExpBFS.cpSrc, Complexity.LogProg.ExpBFS.cp_run, Complexity.LogProg.ExpBFS.d1_run, Complexity.LogProg.ExpBFS.four_mul_kb_le, Complexity.LogProg.ExpBFS.fscan_run, Complexity.LogProg.ExpBFS.fz_run, Complexity.LogProg.ExpBFS.inc_step, Complexity.LogProg.ExpBFS.kb, Complexity.LogProg.ExpBFS.kb_mono, Complexity.LogProg.ExpBFS.lgv, Complexity.LogProg.ExpBFS.lgv_le_succ, Complexity.LogProg.ExpBFS.lt_kb, Complexity.LogProg.ExpBFS.mv_run, Complexity.LogProg.ExpBFS.nz_run, Complexity.LogProg.ExpBFS.orcN, Complexity.LogProg.ExpBFS.reach_cf, Complexity.LogProg.ExpBFS.rr, Complexity.LogProg.ExpBFS.rs1, Complexity.LogProg.ExpBFS.search_run, Complexity.LogProg.ExpBFS.st1, Complexity.LogProg.ExpBFS.st2, Complexity.LogProg.ExpBFS.st3, Complexity.LogProg.ExpBFS.st5, Complexity.LogProg.ExpBFS.stepH_of, Complexity.LogProg.ExpBFS.two_le_codeK, Complexity.LogProg.ExpBFS.wkSt, Complexity.LogProg.ExpBFS.wk_run, Complexity.LogProg.MHalt, Complexity.LogProg.MReach, Complexity.LogProg.MReach.halt, Complexity.LogProg.MReach.refl, Complexity.LogProg.MReach.trans, Complexity.LogProg.astep, Complexity.LogProg.indicator_eq_decide, Turing.FinNDTM.DecidesInSpace, Turing.MultiTapeTM.indicator}
\difficulty{7}
Let $N$ be a finite nondeterministic Turing machine deciding a language $L$ within space
bound $sp$, and let $x$ be a Boolean string. For all natural numbers $d,s$ and all
Boolean memories $M:\mathbb{Z}\to\{\mathrm{false},\mathrm{true}\}$, suppose
$sp(|x|)-s=d$, $s\le sp(|x|)$, and every $z$ with $M(z)=\mathrm{true}$ satisfies
$z<2\,kb(N,|x|,s)$. Then the exponential breadth-first-search program, started with
radius $s$, logarithm register $lgv(|x|)$, all other registers zero, memory $M$, and
memory head at $0$, halts with the correct answer to whether $x\in L$. Every
configuration before halting satisfies the program guard with register bound
$bigB(N,|x|,sp(|x|))$ and memory-head position between $0$ and $3\,kb(N,|x|,sp(|x|)+1)$,
inclusive.
\end{theorem}

\begin{lemma}[Exponential time bound]
\lean{Complexity.time_bound}
\label{Complexity.time_bound}
\leanok
\uses{Complexity.Codec.codeK, Complexity.LogProg.ExpBFS.bigB, Complexity.LogProg.ExpBFS.kb, Complexity.LogProg.ExpBFS.lgv, Complexity.LogProg.ExpBFS.lgv_le, Complexity.LogProg.ExpBFS.logSpace_le_lgv, Complexity.LogProg.ExpBFS.two_le_codeK, Complexity.logSpace, Turing.FinNDTM}
\difficulty{5}
Let $N$ be a finite nondeterministic Turing machine over the binary alphabet, with $k$
work tapes and state set $Q$, and let $C,c_0,n,S$ be natural numbers satisfying
$\lfloor\log_2 n\rfloor+1\le S$. Set $H=3(|Q|+1)$ and $K=2(H+3k+1)+2$. For $m\in\mathbb
N$, define $h(0)=0$, $h(m)=\lfloor\log_2 m\rfloor+1$ when $m>0$, and $g(m)=1+h(m)$; set
$b(m,s)=2^{K(s+g(m)+1)}$ and $B(m,t)=3b(m,t+1)+m+3$. Then
\[
2^{C(\lfloor\log_2(B(n,c_0S)+1)\rfloor+1+\lfloor\log_2
n\rfloor+2)}\bigl(3b(n,c_0S+1)+1\bigr)^2
\le 2^{\left((C(K(c_0+3)+4)+2K(c_0+3)+4)(S+1)\right)}.
\]
\end{lemma}

\begin{theorem}[Nondeterministic space and exponential time]
\lean{Complexity.NSPACE_subset_exp_dtime_of_le}
\label{Complexity.NSPACE_subset_exp_dtime_of_le}
\leanok
\uses{Complexity.Codec.acceptLang_mem, Complexity.Codec.codeK, Complexity.Codec.initLang_mem, Complexity.Codec.stepLang_mem, Complexity.DTIME, Complexity.LogProg.ExpBFS.A, Complexity.LogProg.ExpBFS.St.toV_zero, Complexity.LogProg.ExpBFS.St.zero, Complexity.LogProg.ExpBFS.bigB, Complexity.LogProg.ExpBFS.decs, Complexity.LogProg.ExpBFS.iter_run, Complexity.LogProg.ExpBFS.kb, Complexity.LogProg.ExpBFS.lgv, Complexity.LogProg.ExpBFS.orcN, Complexity.LogProg.ExpBFS.orcN_len, Complexity.LogProg.ExpBFS.reach_cf, Complexity.LogProg.ExpBFS.setup_run, Complexity.LogProg.MHalt.mono, Complexity.LogProg.MReach.halt, Complexity.LogProg.length_bits_le_log, Complexity.LogProg.marm_dtime, Complexity.NSPACE, Complexity.inputLenLang_mem, Complexity.logSpace, Complexity.time_bound, Turing.FinTM.ComputesInTime.mono}
\difficulty{5}
Let $S:\mathbb{N}\to\mathbb{N}$ be a space bound such that $\lfloor\log_2 n\rfloor+1\le
S(n)$ for every $n\in\mathbb{N}$. Then every language decidable by a nondeterministic
Turing machine using space at most a constant multiple of $S$ belongs to
$\mathrm{DTIME}(n\mapsto 2^{c(S(n)+1)})$ for some constant $c\in\mathbb{N}$.
\end{theorem}
